% !PRINTED 601
\weberpage{601}
Nehmen wir in einer Ebene ein rechtwinkeliges Koordinatensystem
 $X, Y$ an, und bezeichnen als \gesperrt{Gitterpunkte} die Punkte
mit den Koordinaten
\[
X = (x_0 + Mx)\, a_{22}\, \sqrt{\dfrac{a}{t}},
\]
\[
Y = (y_0 + My)\, a_{22}\, \sqrt{\dfrac{a}{t}},
\]
worin $x_0, y_0$ festgehalten werden und $x, y$ alle ganzen rationalen
Zahlen durchläuft, so wird durch dieses Gitter die Ebene in
Quadrate geteilt, deren Fläche
\[
\delta^2 = M^2\, \dfrac{a_{22}^2\, a}{t}
\]
ist, und wenn $T_0$ die Anzahl der Gitterpunkte ist, die innerhalb
der Ellipse
\[
\varphi(X, Y) = 1
\]
oder auf ihrer Grenze liegen, so ist $T_0\,\delta^2$ für ein unendliches $t$
gleich dem Flächeninhalt dieser Ellipse, also gleich $\dfrac{\pi}{\sqrt{-\Delta}}$.

Berücksichtigt man noch die Flächenelemente $\delta^2$, die von der
Peripherie der Ellipse durchschnitten werden, aus denen man ein
Maß für den Fehler dieser Flächenbestimmung erhält, so ist
\weq{5}{T_0 = \dfrac{\pi t}{M^2\, a_{22}^2\, a\, \sqrt{-\Delta}} + R_0\sqrt{t},}
worin $R_0$ eine mit unendlich wachsendem $t$ endlich bleibende
Größe ist [Bd.\,II, \S\,194, (6)].

Einen solchen Ausdruck erhalten wir für jedes $\alpha_0$, und wenn
wir diese Ausdrücke alle addieren, so erhalten wir mit Rücksicht
auf (3) den Satz:

\gesperrt{Ist $A$ eine Kongruenzgruppe mit dem Modul $M$
und $\mathfrak{m}$ ein zu $M$ teilerfremdes Ideal in $\Omega$, ferner $T$
die Anzahl der in $\overline{E}A$ enthaltenen ganzen Hauptideale,
deren Norm nicht größer als $t$ ist, so ist}
\weq{6}{T = \dfrac{g t}{N(\mathfrak{m})} + R\sqrt{t},}
\gesperrt{worin $g$ eine von $\mathfrak{m}$ und $t$ unabhängige, von Null
verschiedene Zahl und $R$ eine mit unendlich wachsendem $A$ endliche % [sic]
Funktion von $t$ ist.}
